%1/2 inch left/right margins, 1in top, 1/2 bottom
\documentclass[12pt]{article}
\hbadness=10000
\tolerance=10000
\textheight 9.8in
\textwidth 7.4in
\oddsidemargin -.5in
\topmargin -1in
\headsep 0in  
\headheight 0in
\begin{document}
\begin{flushright}
\textbf{Fall 2016}\\[1.1in]
\end{flushright}

\begin{center}
\textbf{\Huge STATISTICS 801}\\[.1in]
\textbf{\large{Exam 2 (100 points)}}\\[.2in]
\end{center}
\raggedleft Name \rule{3.5in}{.01in}\\[1.1in]
\raggedright  Instructions: 
\begin{itemize}
\item	You are allowed 110 minutes for this exam. Please pace yourself accordingly. 
\item	You may use a calculator and your formula sheet on this exam. However, you are not allowed to use your cell phone AT ALL! It is a "closed book" exam.
\item	The t-table and the Studentized range q-table are attached to your exam so you DON'T need to use your own tables.
\item	In order to receive full credit, you must show your work and carefully justify your answers. The correct answer without any work will receive little or no credit.
\item	If you need more room, use the backs of the pages AND INDICATE that you have done so. 
\item	Please write neatly. 
\item	GOOD LUCK!!
\end{itemize}
\newpage
\setlength{\headsep}{0.2in}
\begin{enumerate}
\item An experiment is conducted to compare a revolutionary
    new method of teaching reading to ``\textit{slow
    learners}'' to the current standard method.  Of a random
    sample of 20 slow learners, 8 are taught by the new
    method and 12 are taught by the standard method.  Both
    groups were taught by qualified instructors and under similar
    conditions for a 6-month period.  The scores obtained by the 
participants in a reading test given at the end of this period, and 
some relevant statistics computed for the two samples,
respectively,  are given below:\\

\vspace*{-.3in}

    \begin{center}
    \begin{tabular}{rrrr|rrrr} \\ \hline
    \multicolumn{4}{c}{New Method} & \multicolumn{4}{c}{Standard Method} \\ \hline
    80 & 80 & 79 & 81 &  79 & 62 & 70 & 68 \\
    76 & 66 & 79 & 76 &  73 & 76 & 86  & 73 \\ 
       &    &    &    &  72 & 68 & 75 & 66 \\ \hline
     \multicolumn{4}{c}{$n = 8$} & \multicolumn{4}{|c}{$n = 12$} \\ 
        \multicolumn{4}{c}{$\bar{y}_1=77.125$} & \multicolumn{4}{|c}{$\bar{y}_2=72.333$}\\
     \multicolumn{4}{c}{$s_1^2=23.55$} & \multicolumn{4}{|c}{$s_2^2=40.24$} \\  \hline
    \end{tabular}
    \end{center}
 Let $\mu_1$ and $\mu_2$ represent
mean reading scores of slow learners taught using the new method 
and the standard method, respectively. Assume that a larger score 
implies an improvement.
  \begin{enumerate}
  \item (10) \parbox[t]{6in}{Is there a reason to believe that
the assumption of equal population variances is justified? Explain.  
Perform a statistical test to verify this assumption, stating the hypotheses and
the rejection region clearly. Use
$\alpha=.05$.}\\
For making your decision use the appropriate $F$-value(s). Circle the one(s) you are using!\\
$F_{0.025,7,11}=3.76$ or $F_{0.025,8,12}=3.51$ or $F_{0.025,1,8}=7.57$ or \\
$F_{0.025,11,7}=4.7$ or $F_{0.025,12,8}=4.2$ or $F_{0.025,8,1}=956.7$.  
  
  
\newpage
\setlength{\headsep}{0.2in}

\item (10) \parbox[t]{6in}{State the appropriate
    hypotheses to determine if the scores indicate that the new
    method is an improvement over the standard method.  Compute the test statistic
    and  an approximate p-value and use it to make your decision.}\\[2.8in]
  \item (10) \parbox[t]{6in}{Construct a 95\% confidence interval
    for $\mu_1 - \mu_2$.  Show work. Explain what the experimenter
    needs to have done to make this estimate more accurate (i.e., to obtain a 
    narrower interval)? Use the the standard error of the difference in means
to explain.}\\[2.8in]

  
  \end{enumerate}
\newpage  
\setlength{\headsep}{0.2in}
 \item A consumer protection magazine was interested in comparing tires
purchased from four different companies that each claimed their tires
would last 40,000 miles. A random sample of 10 tires of each brand
was obtained and tested under simulated road conditions. The number
of miles until the tread thickness reached a specified depth was
recorded for all tires. The data are given (in 1,000 miles) below:
\begin{center}
\begin{tabular}{c|rrrrrrrrrr|rr}
\hline 
 \multicolumn{1}{c|}{\textit{Brand}} & \multicolumn{10}{c|}{\textit{Miles(in thousands)}} & 
     \multicolumn{1}{c}{$\sum{y}$} &
     \multicolumn{1}{c}{$\bar{y}$}  \\ \hline
I  &38.9 &39.7 &42.3 &39.5 &39.6 &35.6 &36.0 &39.2 &37.6 &39.5 & 387.9 & 38.79\\
II &44.6 &46.9 &48.7 &41.5 &37.5 &33.1 &43.4 &36.5 &32.5 &42.0 & 406.7 & 40.67 \\ 
III &51.3 &53.6 &58.2 &52.1 &51.2 &53.1 &51.8 &55.8 &52.1 &51.4 & 530.6 & 53.06 \\ 
IV &39.1 &38.7 &40.4 &41.9 &39.4 &36.1 &40.5 &35.2 &32.3 &36.0 & 379.6 & 37.96 \\ \hline
\end{tabular}
\end{center}


Part of the ANOVA table is shown below:

\begin{center}
\begin{tabular}{|c|c|c|c|c|c|}
\hline 
Source & \,\,\,\,\,\,df \,\,\,\,\,\,& SS & \,\,\,\,\,\,\,\,MS\,\,\,\,\,\,\,\, &\,\,\,\,\,\,\,\, F\,\,\,\,\,\,\,\, & p-value\\ \hline
Treatment &  & 1491.81 & & & $1.188\times 10^{-11}$\\ \hline
Error &&440.28&&---&---\\ \hline
Total &&&---&---&---\\ \hline
\end{tabular}
\end{center}
\begin{enumerate}
\item (10) \parbox[t]{6in}{Fill in the missing values in the ANOVA table (include both equations and actual number).}
\item (10) \parbox[t]{6in}{Write out the linear statistical model and describe it, including \textbf{all} assumptions. Be specific! Use mathematical notations AND word explanations too.}\\[2.8in]
\item (3) \parbox[t]{6in}{What type of experimental design is it? How do you know? Justify your answer.}\\[2.8in]
\item (2) \parbox[t]{6in}{What hypothesis can you test with the F and p-values found in the ANOVA table? Write down the null and alternative hypotheses.}\\[2.8in]
\item (5) \parbox[t]{6in}{Using the p-value given in the ANOVA table, state your decision and conclusion in terms of the problem. $\alpha=0.05$.} \\[2.8in]
\item (3) \parbox[t]{6in}{The company producing $Brand I$ is located in Nebraska. We are interested in comparing the Nebraska tires with all of the other tires. Write a single linear contrast that describes the comparison described above.} \\[1.8in]
\item (3) \parbox[t]{6in}{Show that the linear combination you specified in part (f) is a linear contrast.}\\[2.8in]
\item (10) \parbox[t]{6in}{Using the t-distribution, test whether the linear contrast described in part (f) is significant. Make sure you include every step of the test.$\alpha=0.05$.}\\[3.6in]
\item (10) \parbox[t]{6in}{Perform Tukey's procedure to perform every pairwise comparison.$\alpha=0.05$. Give a conclusion.} 
\end{enumerate}
\newpage
\item \textbf{TRUE/FALSE}. Circle your answer. (2 points each.)
\begin{enumerate}
\item \textbf{T / F} Sampling error is the only cause two sample means are different.
\item \textbf{T / F} The number of mutually orthogonal contrasts is the same as the number of means.
\item \textbf{T / F} A significant F-test from a one-way ANOVA tells us that there is a significant difference somewhere among the means, but it doesn't tell us where it is.
\item \textbf{T / F} Within-group variance is caused by sampling error and treatment effects.
\item \textbf{T / F} The F statistic in the ANOVA table compares the ratio of between-group variance to within-group variance.
\item \textbf{T / F} The paired t-test assumes that the data in each group are normally distributed.
\item \textbf{T / F} The paired t-test is appropriate for comparing the means of independent groups of observations.
\end{enumerate}
\end{enumerate}
\end{document}